\documentclass[a4paper,11pt]{article}
\usepackage[utf8]{inputenc}
\usepackage[spanish]{babel}
\usepackage{amsmath, amsthm, amssymb, amsfonts}
\usepackage[dvipsnames]{xcolor}
\usepackage{tcolorbox}
\tcbuselibrary{skins, breakable, theorems}
\usepackage[margin=2.5cm]{geometry}

% Definición de colores
\definecolor{maincolor}{RGB}{41, 128, 185}
\definecolor{defcolor}{RGB}{39, 174, 96}
\definecolor{thmcolor}{RGB}{142, 68, 173}
\definecolor{notecolor}{RGB}{230, 126, 34}

% Configuración de las cajas
\newtcbtheorem[number within=section]{definicion}{}%
{colback=defcolor!5,colframe=defcolor!80!black,fonttitle=\bfseries, breakable}{def}

\newtcbtheorem[number within=section]{teorema}{Teorema}%
{colback=thmcolor!5,colframe=thmcolor!80!black,fonttitle=\bfseries, breakable}{thm}

\newtcolorbox{nota}[1][]{colback=notecolor!5,colframe=notecolor!80!black,fonttitle=\bfseries, breakable, title=#1}

\title{\textbf{\color{maincolor}Resumen I1 - Inferencia Estadística}}
\author{Santiago Franulic}
\date{}

% Configuraciones extra
\renewcommand{\arraystretch}{1.4}
\usepackage{hyperref}
\hypersetup{
    colorlinks=true,
    linkcolor=maincolor,
    urlcolor=maincolor,
    citecolor=maincolor
}

\begin{document}
\maketitle
\tableofcontents
\newpage

\section{Estudios Estadísticos}

\begin{definicion}{Estudio Estadístico}{}
    Proceso de recolección, análisis e interpretación de datos.
\end{definicion}

\begin{definicion}{Tipos de Estudios Estadísticos}{}
    \begin{itemize}
        \item Segun el objetivo:\\
    \begin{tabular}{l|l}
        \hline
        \textbf{Analitico} & Busca explicar una relación causa-efecto entre variables.\\
        \hline
        \textbf{Descriptivo} & Describe una población o fenómeno.\\
        \hline
    \end{tabular}
    \vspace{0.1cm}
    \item Segun el periodo de tiempo:\\
    \begin{tabular}{l|l}
        \hline
        \textbf{Transversal} & Analiza en un unico momento en el tiempo.\\
        \hline
        \textbf{Longitudinal} & Analiza en distintos momentos en el tiempo.\\
        \hline
    \end{tabular}
    \vspace{0.1cm}
    \item Segun el control de variables:\\
    \begin{tabular}{l|l}
        \hline
        \textbf{Observacional} & No interviene el investigador.\\
        \hline
        \textbf{Experimental} & Interviene el investigador.\\
        \hline
    \end{tabular}
    \vspace{0.1cm}
    \item Segun la dirección temporal:\\
    \begin{tabular}{l|l}
        \hline
        \textbf{Prospectivo} & Analiza el fenómeno en el futuro.\\
        \hline
        \textbf{Retrospectivo} & Analiza el fenómeno en el pasado.\\
        \hline
    \end{tabular}
    \end{itemize}
\end{definicion}

\begin{definicion}{Recopilación de datos}{}
    \begin{itemize}
        \item Muestreo $\rightarrow$ Busca describir una poblacion.
        \begin{itemize}
            \item Probabilistico $\rightarrow$ Todos tienen una probabilidad de ser elegidos.
            \begin{itemize}
                \item Estratificado: Se divide la poblacion en estratos y se muestrea aleatoriamente de cada estrato.
                \item Conglomerados
                \item Multietápico
            \end{itemize}
            \item No probabilistico $\rightarrow$ No todos tienen una probabilidad de ser elegidos.
        \end{itemize}
        \item Reclutamiento $\rightarrow$ Busca explicar una relación causa-efecto entre variables.
    \end{itemize}
\end{definicion}

\end{document}

